\documentclass[11pt]{article}

\usepackage[margin=1in]{geometry}
\usepackage{amsmath,amssymb,amsthm,mathtools}
\usepackage{booktabs}
\usepackage{microtype}
\usepackage{enumitem}
\usepackage{hyperref}

\hypersetup{colorlinks=true,linkcolor=blue,citecolor=blue,urlcolor=blue}

\newtheorem{theorem}{Theorem}[section]
\newtheorem{lemma}[theorem]{Lemma}
\newtheorem{corollary}[theorem]{Corollary}
\theoremstyle{remark}
\newtheorem{remark}[theorem]{Remark}

\newcommand{\Qx}{\mathcal Q_x}
\newcommand{\Dx}{\mathcal D_x}
\newcommand{\one}{\mathbf 1}
\newcommand{\ee}{\mathrm e}
\newcommand{\Kvec}{\mathbf K}
\newcommand{\Kj}{K_j}
\newcommand{\Kjq}{K_{j,q}}
\newcommand{\Kell}{K_\ell}
\newcommand{\norm}[1]{\left\lVert #1\right\rVert}

\title{\textbf{Prime-Shell Hilbert Lift and the Sharp Collapse Barrier}}
\author{Liang Wang\\
Huazhong University of Science and Technology\\
Wuhan 430074, P.R. China}
\date{August 2026}

\begin{document}
\maketitle

\begin{abstract}
The previous structural stage attached a literal common-source twin-prime
kernel to a finite interval by a reduced rational-frequency large sieve, but
it collapsed the prime shell before the finite-window estimate.  We retain the
prime label and the four-packet label as Hilbert coordinates.  For
$H=x^{21/32}$, $Q=x^{1/3}$, and $U=x^{133/400}$, fixed-$q$ cutoff injectivity
and an elementary unsigned cluster bound give
\[
 \frac{1}{N}\sum_{n\in I_x}\norm{\Kvec(n)}_2^2
 \ll J M^2 x^{1/96}(\log x)^5.
\]
Here $\Kvec$ has coordinates $(j,q)$, $J$ is the fixed packet count, and
$M$ is a uniform profile bound.  Recombining the prime labels costs exactly
the elementary factor $P=\#\Qx\leq 2Q$ and recovers the earlier scalar power
$x^{11/32}$.  A finite aligned fixture attains the full ratio $P=4$, while a
parallel four-packet algebraic fixture has unit-projection ratio one.  Thus
the new theorem preserves the interfaces and isolates the precise prime-shell
collapse barrier; it supplies no arithmetic cancellation, $L^2$ advance,
fixed-atom credit, or twin-prime conclusion.
\end{abstract}

\paragraph{Claim level.}
\texttt{PROVED\_STRUCTURAL\_L1 / PRIME\_LABEL\_AND\_PACKET\_PRESERVING\_LIFT}.\\
The Hilbert lift, exponent ledger, and scalar recovery are proved under the
declared common-source hypotheses.  The finite fixtures are exact structural
controls, not asymptotic arithmetic evidence.

\section{Introduction}

A useful analytic route toward the twin-prime problem should expose every
place where a collective estimate is required.  The recent finite-window
attachment did this for a common-source kernel by regrouping divisor rows
according to their reduced rational frequencies and applying the additive
large sieve.  That argument was intentionally unsigned.  In particular, it
collapsed the prime shell before invoking the finite-window inequality.

The present paper asks a narrower follow-up question: can the outer prime label
and the four-packet label be retained in the same finite-window object?  This
is important for bookkeeping.  If the labels are removed too early, a later
signed prime-shell or polarized reassembly can be accidentally replaced by a
Cauchy estimate without recording its cost.  If they are retained, one can
measure exactly what the structural argument does and does not provide.

We work with the same V46 common-source scales
\[
 H=x^{21/32},\qquad Q=x^{1/3},\qquad
 Y_0=\frac{H}{4Q},\qquad U=x^{133/400}.
\]
The prime shell is
\[
 \Qx=\{q\ {\rm prime}: Q<q\leq 2Q\},
 \qquad P=\#\Qx,
\]
and the squarefree transition band is
\[
 \Dx=\{d\in\mathbb Z:Y_0<d\leq U,\ \mu(d)^2=1\}.
\]
The physical interval is $I_x=(x/2,x]\cap\mathbb Z$, with $N=|I_x|$.

Our main result is a tensorized version of the standard additive large sieve.
It preserves $(q,j)$ as coordinates and gives a normalized split exponent
$1/96$, since
\[
 \frac{Q^2}{H}=x^{1/96}.
\]
Only after this result do we form the scalar packet shell.  Pointwise Cauchy
over $q$ costs $P\leq 2Q$, changing $Q^2/H$ into $Q^3/H=x^{11/32}$.
This identity is the central accounting statement of the paper.

The contributions are:
\begin{enumerate}[leftmargin=2em]
  \item an exact prime-label- and packet-label-preserving finite-window
        Hilbert lift;
  \item an explicit split-versus-scalar exponent ledger, with the $P$ cost
        isolated rather than hidden;
  \item a positive-semidefinite packet Gram formulation compatible with exact
        four-point complex polarization;
  \item two scoped adversarial controls: a $P$-saturating prime alignment and
        a geometry-only packet alignment.
\end{enumerate}

\paragraph{Status boundary.}
Nothing here proves a signed estimate for the literal prime shell.  In the
terminology of the route ledger,
\[
 \texttt{ARITHMETIC\_ADVANCE=NO},\qquad
 \texttt{FIXED\_ATOM\_CREDIT=0},\qquad
 \texttt{L2=NONE}.
\]
The result is a structural bridge to the next open theorem.

\section{The split common-source object}

Let $\psi_j:\mathbb R\to\mathbb C$ be bounded profiles for
$0\leq j<J$, and put
\[
 M=\max_{0\leq j<J}\norm{\psi_j}_\infty.
\]
For $q\in\Qx$, $h\leq U$, and $a\pmod h$, define
\begin{equation}
 B_{h,q}^{(j)}(a)
 =
 \sum_{0<|m|\leq \lfloor hq/H\rfloor}
 \psi_j\!\left(\frac{Hm}{hq}\right)
 \one_{m\overline q\equiv a\pmod h}.
 \label{eq:row}
\end{equation}
The row is defined on all residues; only primitive residues will be used as
reduced-frequency numerators.  Set
\begin{equation}
 c_d=\frac{\mu(d)\log d}{d},\qquad
 C_h=\sum_{\substack{d\in\Dx\\h\mid d}}c_d.
 \label{eq:coeff}
\end{equation}
The split coordinates are
\begin{equation}
 \Kjq(n)=
 \sum_{h\leq U}\ \sum_{\substack{a\bmod h\\(a,h)=1}}
 C_h B_{h,q}^{(j)}(a)\ee(na/h),
 \qquad
 \Kvec(n)=\bigl(\Kjq(n)\bigr)_{\substack{0\leq j<J\\q\in\Qx}}.
 \label{eq:split}
\end{equation}
The scalar packet shell is
\begin{equation}
 \Kj(n)=\sum_{q\in\Qx}\Kjq(n).
 \label{eq:scalar}
\end{equation}

The exponent relations needed below are
\begin{equation}
 \frac{Q^2}{H}=x^{1/96},\quad
 \frac{Q^3}{H}=x^{11/32},\quad
 \frac{U^2}{x}=x^{-67/200},\quad
 \frac{UQ}{H}=x^{23/2400}.
 \label{eq:exponents}
\end{equation}
For sufficiently large $x$, $4Q<H$ and $U<Q$.  In particular every shell
prime is a unit modulo every divisor in $\Dx$.

\subsection{The inherited reduced-frequency identity}

The divisor-dilation identity from the preceding finite-window stage is
\begin{equation}
 B_{d,q}^{(j)}(ka)=B_{h,q}^{(j)}(a)
 \qquad (d=kh,\ h\mid d),
 \label{eq:dilation}
\end{equation}
because the congruence forces $m=kn$ and the cutoff and profile argument
scale exactly.  Thus the literal common-source kernel can be partitioned by
the reduced frequency $a/h$ while retaining $(q,j)$.  This is an identity,
not an orthogonality assertion.

\section{The Hilbert-valued finite-window theorem}

\begin{lemma}[fixed-$q$ cutoff injectivity]
\label{lem:inject}
For sufficiently large $x$,
\[
 \sum_{\substack{a\bmod h\\(a,h)=1}}
   \left|B_{h,q}^{(j)}(a)\right|^2
 \leq 2M^2\frac{hq}{H}.
\]
\end{lemma}

\begin{proof}
If two nonzero atoms $m_1,m_2$ in \eqref{eq:row} occupy the same residue,
then $h\mid(m_1-m_2)$.  Since $q\leq 2Q$ and $4Q<H$,
\[
 |m_1-m_2|
 \leq 2\left\lfloor\frac{hq}{H}\right\rfloor<h.
\]
Thus the atoms are distinct modulo $h$.  Summing their squared magnitudes over
all residues, and then restricting to primitive residues, gives at most
$2\lfloor hq/H\rfloor M^2\leq 2M^2hq/H$.
\end{proof}

\begin{lemma}[active-cluster harmonic bound]
\label{lem:cluster}
The literal coefficients satisfy
\[
 \sum_{h\leq U}h|C_h|^2\ll(\log x)^5.
 \label{eq:cluster}
\]
\end{lemma}

\begin{proof}
If the row indexed by $h$ is nonzero for at least one shell prime, then
$\lfloor hq/H\rfloor\geq1$, hence $h\geq H/(2Q)$.  Writing $d=hk$ in
\eqref{eq:coeff}, we have $k\leq U/h\leq 2UQ/H$.  Therefore
\[
 |C_h|
 \leq \frac{\log U}{h}\sum_{k\leq U/h}\frac1k
 \ll \frac{(\log x)^2}{h}.
\]
The active $h$ form a subinterval of $[H/(2Q),U]$, so
\[
 \sum_hh|C_h|^2
 \ll(\log x)^4\sum_{H/(2Q)\leq h\leq U}\frac1h
 \ll(\log x)^5.
\]
No cancellation in $\mu$ is used.
\end{proof}

\begin{lemma}[Hilbert-valued additive large sieve]
\label{lem:hilbert}
Let $I$ be a consecutive interval of $N$ integers.  Then
\[
 \sum_{n\in I}\norm{\Kvec(n)}_2^2
 \leq (N+U^2)
 \sum_{\substack{h\leq U\\(a,h)=1}}
 \sum_{q\in\Qx}\sum_{0\leq j<J}
 |C_hB_{h,q}^{(j)}(a)|^2.
 \label{eq:hilbert}
\]
\end{lemma}

\begin{proof}
Distinct reduced fractions with denominators at most $U$ have circular
spacing at least $U^{-2}$.  Apply the standard additive large-sieve inequality
\cite{MontgomeryVaughan1973} to the scalar expansion of each coordinate
$\Kjq(n)$, with coefficient $C_hB_{h,q}^{(j)}(a)$, and sum the resulting
inequalities over $(q,j)$.  This is precisely the coordinate realization of
the tensor lift from scalar coefficients to
$\ell^2(\Qx\times\{0,\ldots,J-1\})$.
\end{proof}

\begin{theorem}[prime-label- and packet-preserving lift]
\label{thm:split}
For fixed $J$ and bounded packet profiles,
\[
 \sum_{n\in I_x}\norm{\Kvec(n)}_2^2
 \ll J M^2(N+U^2)\frac{Q^2}{H}(\log x)^5.
 \label{eq:split-bound}
\]
Consequently,
\[
 \frac1N\sum_{n\in I_x}\norm{\Kvec(n)}_2^2
 \ll J M^2x^{1/96}(\log x)^5,
 \label{eq:split-normalized}
\]
and the unnormalized exponent is $97/96+o(1)$.
\end{theorem}

\begin{proof}
By Lemma~\ref{lem:inject},
\[
\begin{aligned}
 &\sum_{h,a,q,j}|C_hB_{h,q}^{(j)}(a)|^2\\
 &\quad\leq
 \frac{2M^2J}{H}\left(\sum_{q\in\Qx}q\right)
 \sum_hh|C_h|^2\\
 &\quad\ll J M^2\frac{P Q}{H}(\log x)^5
 \ll J M^2\frac{Q^2}{H}(\log x)^5,
\end{aligned}
\]
where $P\leq2Q$.  Insert this into Lemma~\ref{lem:hilbert}.  Since
$U^2/N=x^{-67/200+o(1)}$, division by $N$ gives the normalized estimate.
\end{proof}

\subsection{Scalar recovery and the packet Gram}

\begin{corollary}[exact scalar collapse ledger]
\label{cor:scalar}
For the packet shell \eqref{eq:scalar},
\begin{equation}
 \frac1N\sum_{n\in I_x}\sum_{j<J}|\Kj(n)|^2
 \ll J M^2x^{11/32}(\log x)^5.
 \label{eq:scalar-bound}
\end{equation}
\end{corollary}

\begin{proof}
For each $j,n$, Cauchy gives
$|\sum_q\Kjq(n)|^2\leq P\sum_q|\Kjq(n)|^2$.  Sum over $j,n$ and
apply Theorem~\ref{thm:split}; then use $P\leq2Q$ and
$Q^3/H=x^{11/32}$.
\end{proof}

For a finite interval define the packet Gram matrix
\begin{equation}
 G_I[j,\ell]=\sum_{n\in I}\Kj(n)\overline{\Kell(n)}.
\label{eq:gram-definition}
\end{equation}
It is positive semidefinite, and Corollary~\ref{cor:scalar} is a trace bound.
For every unit packet vector $\omega$,
\begin{equation}
 \sum_{n\in I}\left|\sum_j\overline{\omega_j}\Kj(n)\right|^2
 =\omega^*G_I\omega\leq\operatorname{tr}(G_I).
 \label{eq:gram}
\end{equation}
The exact four-packet polarization interface remains available through
\begin{equation}
 x\overline y=\frac14\sum_{r=0}^{3}i^r|x+i^ry|^2.
 \label{eq:polarization}
\end{equation}
The trace inequality is unsigned and therefore does not supply the signed
four-packet estimate needed by the terminal route.

\section{Adversarial controls}

\subsection{A P-saturating prime-label fixture}

Take
\[
 d=5,\qquad H=500,\qquad
 \{q\}=\{101,131,151,181\},\qquad \psi(t)\equiv1.
\]
All four primes are $1\pmod5$, and $\lfloor dq/H\rfloor=1$.  Hence every
fixed-$q$ row is $e_1+e_4$.  The coherent row and diagonal energies are
\[
 \left\|\sum_qB_{5,q}\right\|_2^2=32,\qquad
 \sum_q\|B_{5,q}\|_2^2=8,
\]
so their ratio is $4=P$.  This is an exact rational finite structural
adversary.  It refutes only the scoped shortcut that fixed-$q$ row geometry
would imply a free or sub-$P$ shell collapse.

\subsection{A packet-projection alignment}

Let $v\neq0$ and let $\omega$ be a unit four-vector.  Set
$Z^{(j)}=\omega_jv$.  Then
\[
 \sum_j\norm{Z^{(j)}}^2=\norm v^2,\qquad
 \left\|\sum_j\overline{\omega_j}Z^{(j)}\right\|^2=\norm v^2.
\]
For $\omega=(1,i,-1,-i)/2$, the projection-to-total ratio is one.  This is an
algebraic Hilbert-space obstruction: packet geometry alone cannot create
cancellation.  It is not a claim that every literal TPC profile realizes this
parallel family.

\begin{table}[t]
\centering
\caption{TPC-218 claim firewall.}
\begin{tabular}{ll}
\toprule
Structural lift & \texttt{PROVED\_STANDARD\_TENSOR\_LIFT}\\
Split exponent & \texttt{PROVED\_1\_OVER\_96\_LOG\_FIVE}\\
Scalar recovery & \texttt{PROVED\_P\_FACTOR\_RECOVERS\_11\_OVER\_32}\\
Prime-label orthogonality & \texttt{REFUTED\_SCOPED}\\
Packet cancellation & \texttt{NONE}\\
Arithmetic advance & \texttt{NO}\\
Prime shell / four packets & \texttt{OPEN}\\
Full Gate B / strict one-over-400 & \texttt{OPEN / UNPAID}\\
\bottomrule
\end{tabular}
\end{table}

\section{Route position and reproducibility}

The result belongs to Route B's structural threshold-A layer.  Route A is not
applicable.  The strongest positive result is the label-preserving split
envelope with exponent $1/96$, together with the exact identification of the
single $P$ factor needed for scalar recovery.  The strongest obstruction is
the finite $P$-saturating row alignment.  The open theorem is a signed
prime-shell reassembly that beats this cost while retaining the literal
four-packet, zero/nonunit, fixed-atom, and normalization interfaces.

The repository release contains a producer, an independent checker, normal and
optimized-mode checks, and a separate adversarial script.  The finite
dilation fixture uses the exact rational surrogate $\mu(d)/d$ solely to check
the index map; the asymptotic theorem above uses the literal
$\mu(d)\log(d)/d$.  No numerical fixture is used as asymptotic evidence.

\section{Conclusion}

TPC-218 turns the instruction \textquotedblleft keep the shell\textquotedblright
into a precise Hilbert-valued theorem.  The split object is controlled at the
$Q^2/H=x^{1/96}$ scale, while scalar shell recovery necessarily records a
$P\leq2Q$ payment and returns the TPC-217 exponent $11/32$.  The finite
aligned row and the parallel packet construction show why geometry alone
cannot pay the missing arithmetic saving.  The next research target is
therefore unambiguous:
\[
 \text{prove a signed prime-shell reassembly beyond the exact \(P\) collapse.}
\]
This paper makes no claim that the target has been achieved.

\paragraph{Acknowledgement of scope.}
The claims are limited to the displayed common-source object and its declared
finite-window normalization.  In particular,
\texttt{TPC218\_ARITHMETIC\_ADVANCE=NO},
\texttt{TPC218\_FIXED\_ATOM\_CREDIT=0}, and
\texttt{TPC218\_FULL\_GATE\_B=OPEN}.

\bibliographystyle{plain}
\bibliography{references}

\end{document}
